\bookchapter{الفصل 0: مرحباً، قبل أن نبدأ}

قبل أن ننطلق في الاستكشاف الرسمي، لنتحدث عن بعض القواعد الأساسية لضمان فهمك لكل ما سيُقال لاحقاً. لا تتعجل، فإعداد الفأس لا يضيع وقت قطع الحطب.

\booksection{أولاً: ماذا يمكن أن يساعدك هذا الكتاب؟ (وما لا يستطيع فعله)}

يمكنه أن يساعدك بشكل منهجي على فهم كيفية استخدام الهاتف، وإدراك المنطق وراء الوظائف، وبغض النظر عما إذا كنت من مستخدمي Android أو Apple، فإن معظم الأفكار تنطبق على الجميع. سيرافقك في التمارين العملية لترتقي من “معرفة الاستخدام” إلى “الفهم العميق للاستخدام”.

إنه ليس دليل استخدام لهاتف محدد بعينه، وليس أخباراً تقنية تتبع المنتجات الجديدة، وليس دليلاً لتعليمك إدمان الهاتف. بصراحة، إذا كنت تتوقع “كيف تستخدم الهاتف لمشاهدة مقاطع الفيديو بشكل أفضل”، فقد تخيب ظنونك — ما يهمنا أكثر هو كيف يجعل الهاتف يخدمك حقاً، لا العكس.

\booksection{ثانياً: شرح “الرموز” لدينا}

لتوضيح الأمر بشكل أفضل، يستخدم الكتاب بعض العلامات الصغيرة، وهذه العلامات مثل الرموز المشتركة بيننا:

\begin{itemize}
  \item \textbf{[اسم الزر]}: يمثل زراً محدداً على الشاشة أو اسم تطبيق، مثل [الإعدادات] و[موافق].
  \item \textbf{>}: يمثل تسلسل العمليات، مثلاً "اضغط [الإعدادات] > [الصوت]" يعني اضغط الإعدادات أولاً، ثم الصوت.
  \item \textbf{الخط العريض}: مهم جداً! قد يكون مفهوماً رئيسياً أو تحذيراً أمنياً، يرجى الانتباه إليه.
  \item \emph{الخط المائل}: أفكار أو تشبيهات تكميلية تساعدك على الفهم بشكل أفضل.
\end{itemize}

\booksection{ثالثاً: Android و iOS، نتحدث عن كليهما}

خمن ماذا، بعض الأشياء تختلف في تسميتها لكنها تؤدي الوظيفة نفسها. مثلاً: تطبيق (Android) و App (iOS)، الإعدادات (كلاهما يسمى الإعدادات)، إدارة الملفات (Android) والملفات (iOS). المهم أن تعرف أنها نفس الشيء، لا تدع المصطلحات تربكك.

بالإضافة إلى ذلك، يوجد الآن نظام HarmonyOS من Huawei في معسكر Android، وله خصائصه الخاصة، وسنتحدث عنه بالتفصيل لاحقاً. باختصار، مهما كان الهاتف الذي تستخدمه، يمكن لهذا الكتاب أن يساعدك.

\booksection{رابعاً: سعة الهاتف: لماذا بعد شراء 128GB، المساحة المتاحة أقل من 110GB؟}

وحدات السعة معروفة: 1TB = 1024GB، 1GB = 1024MB، 1MB = 1024KB... لكن هناك فخ هنا: الشركات المصنعة للأقراص عادةً تحسب بالنظام العشري (1000)، بينما يحسب النظام بالنظام الثنائي (1024)، لذلك القرص “128GB” الذي تشتريه ستكون سعته الفعلية “أقل”.

لكن أسباب نقصان سعة هاتفك الرئيسية ثلاثة:

\begin{enumerate}
  \item \textbf{النظام يحتل مساحة بنفسه}: نظام التشغيل مثل الجدران الحاملة في المنزل، لا بد من وجوده. الأنظمة الحالية تحتل عادةً 10-20GB، وهذا طبيعي.
  \item \textbf{التطبيقات المثبتة مسبقاً من الشركة المصنعة}: التطبيقات الموجودة في الهاتف الجديد تحتل مساحة أيضاً، بعضها يمكن حذفه وبعضها لا.
  \item \textbf{خسائر التقسيم والتهيئة}: لأسباب تقنية، تماماً كما أن بطانة الحقيبة تشغل حجماً.
\end{enumerate}

\begin{tcbinfo}{نصيحة صغيرة عند شراء الهاتف}
هواتف 128GB لم تعد تكفي عملياً، و 256GB كحد أدنى أكثر أماناً. إذا كنت تحب التصوير أو لعب الألعاب، فننصح بالذهاب مباشرة إلى 512GB أو أكبر. تذكر، بالنسبة لمساحة التخزين، من الأفضل أن تكون زائدة عن الحاجة لا ناقصة.

\end{tcbinfo}

\begin{tcbpractice}{لنبدأ عملياً: شاهد "بطاقة الهوية" لهاتفك}
افتح [الإعدادات]، ابحث عن “حول الهاتف” (أو خيار مشابه)، وسجّل: طراز الهاتف، إصدار النظام، إجمالي مساحة التخزين والمساحة المتبقية. هذه هي الخطوة الأولى لمعرفته. بالمناسبة، تحقق من صحة البطارية لتكون على دراية.

\end{tcbpractice}

\booksection{خامساً: هاتفك “مخزن للأصول الرقمية”}

قبل أن نبدأ بالمعرفة المحددة، لنؤسس إدراكاً أساسياً:\textbf{لم يعد الهاتف مجرد أداة اتصال، بل هو أهم مخزن لأصولنا الرقمية}. جهات الاتصال، سجلات الدردشة، صور العائلة، مستندات العمل، تطبيقات البنوك، رموز الدفع، بيانات الصحة، الوثائق الإلكترونية... هذه البيانات مجتمعة تتجاوز قيمتها غالباً قيمة الهاتف نفسه.

\begin{tcbinfo}{لماذا هذا مهم؟}
إذا فُقد هاتفك أو تعطل، فإنك لا تخسر فقط جهازاً بقيمة بضعة آلاف من اليوان،\textbf{بل ربما بيانات لا يمكن تعويضها تراكمت عبر سنوات}— صور نمو أطفالك، تسجيلات صوتية لأقارب متوفين، عقود مهمة لم يتم نسخها احتياطياً. التنظيم العلمي وعادات الحماية يمكن أن تحميك من ثلاثة أبعاد:

\begin{itemize}
  \item \textbf{تحسين كفاءة التشغيل}: التصنيف المنظم للملفات يجعل بحث النظام أسرع، والهاتف لا يعلق لفترات طويلة.
  \item \textbf{إطالة العمر الافتراضي}: إدارة التخزين المعقولة تقلل من عبء القراءة والكتابة على الشريحة، والبطارية والتخزين يدومان لفترة أطول.
  \item \textbf{ضمان أمان البيانات}: الإدارة الهرمية + النسخ الاحتياطي الدوري + تقليص الأذونات، تُبعد عنك "كوارث البيانات".
\end{itemize}

\end{tcbinfo}

\booksubsection{فلسفتنا الأساسية: “الوقاية أولاً، والجمع بين الوقاية والعلاج”}

كما أن العناية بالصحة لا تنتظر المرض لزيارة الطبيب،\textbf{يجب أن تبدأ إدارة الهاتف من “العادات اليومية”}. ما يؤكد عليه هذا الكتاب مراراً ليس مهارة متقدمة معينة، بل هذه المبادئ البسيطة:

\begin{itemize}
  \item \textbf{الوقاية أولاً}: تثبيت التطبيقات من المتاجر الرسمية فقط، إزالة البيانات الحساسة من الصور قبل مشاركتها، عدم تكرار كلمات المرور — هذه "الإجراءات الصغيرة" منخفضة التكلفة لكنها تصدّ معظم المخاطر.
  \item \textbf{فحص دوري}: تنظيف مجلد التنزيلات شهرياً، تدقيق الأذونات ربع سنوياً، نسخ احتياطي كامل نصف سنوياً — الإيقاع الثابت أكثر فاعلية بمئة مرة من إطفاء الحرائق المؤقت.
  \item \textbf{اترك لنفسك مخرجاً}: تفعيل المزامنة السحابية، تسجيل أرقام الحسابات وكلمات المرور جيداً، كتابة "قائمة التركة الرقمية" — لضمان قدرتك أنت وعائلتك على التعامل بهدوء عند وقوع حادث.
\end{itemize}
